\chapter*{Как пользоваться учебником}
\addcontentsline{toc}{chapter}{Как пользоваться учебником}
\markboth{Как пользоваться учебником}{Как пользоваться учебником}

Эта книга учит описывать шоу дронов программой на Python и проверять его в симуляторе. Дроны для этого не нужны: всё происходит на экране. Когда шоу готово и проверено, его можно загрузить в настоящие дроны LiteBee Wing (об этом последний урок).

\textbf{Как устроен каждый урок.} Он начинается с самого простого. Сначала небольшой рабочий пример, потом объяснение, что в нём происходит, и в конце рамка \textbf{«Попробуйте сами»}: несколько упражнений, чтобы закрепить. Ответы ко всем упражнениям в конце книги. Сначала попробуйте сами, потом подсматривайте.

\textbf{Как работать.}
\begin{enumerate}
  \item Возьмите пример из урока. Каждый листинг лежит в папке \texttt{docs/uchebnik/solutions} под тем именем, которое написано над ним.
  \item Откройте файл в Litebee IDE и нажмите \knopka{Проверить}, потом \knopka{Играть}. Посмотрите, что получилось.
  \item Поменяйте число или строку, нажмите \knopka{Проверить} ещё раз. Так учишься быстрее всего.
  \item Выполните упражнения, сверьтесь с ответами.
\end{enumerate}

\textbf{Что проверено.} Все примеры и ответы к упражнениям (файлы \texttt{l*.py} и \texttt{u*.py}) прошли проверки симулятора без предупреждений и ошибок. Файлы, которые нарочно ошибаются, лежат в \texttt{solutions/bad\_examples}: на них показаны настоящие сообщения.

\begin{zametka}[Значки в тексте]
В такой рамке лежит то, что стоит запомнить.
\end{zametka}
\begin{vnimanie}[Осторожно]
Так отмечено то, что легко сделать неправильно.
\end{vnimanie}
\begin{opasno}[Безопасность]
Так отмечено то, что касается настоящих дронов.
\end{opasno}

\textbf{Что ещё не проверено на живом дроне.} Шаги \kod{circle} и \kod{code_z} приложение упаковывает для загрузки так же, как LiteBee Go, но на настоящем дроне они пока не летали. В симуляторе они работают. Шаг \kod{circle_z} в дрон загрузить нельзя вообще. Эти места помечены в тексте.
